\chapter{Executive Summary}
\label{ch:exec-summary}

\noindent\draftmark\medskip

MAITRIA stands for \emph{Mathematical Assurance and Interaction Toolkit
for Reliable Intelligent Agents}: a bundle of public-good software and
mathematics built to discharge the Safeguarded AI programme's TA1.1
and TA1.2 objectives --- a professional-grade implementation of the
modelling-language backend, with machine-checkable assurance carried
through every layer. This book is that bundle's manual: quick starts,
tutorials, and normative references for its three components, in one
volume, written for readers who know the programme's founding
documents and know
category theory, CRDTs, or verification at research depth. Since it
is addressed to machine as well as to human readers, the \LaTeX{}
source is published beside the PDF.

\section*{The three components}

\textbf{Geolog} is the modelling substrate: a relational backend
grounded in geometric logic, in which theories, models, and the
toolkit's own rule systems are object-level data under a single
meta-schema, with the chase as its engine. Persistence, version control, and
synchronization are native to the data model rather than bolted on as
diff-and-merge over text, and capability-based authorization is
designed to converge with the Keyhive line of work. The public cut,
\texttt{geolog-alpha}, has been factored for its debut: a tidy, easily
auditable core that happens to contain platform-optimized decision
procedures for the ground relations the logics need, gated behind
feature flags, while the logics' \emph{rules} live as object-level
entries with respect to geolog's meta-schema, not as engine code.

\textbf{QTSL}, the Quantitative Temporal Specification Logic, is the
specification layer: quantitative temporal properties --- safety,
reach-avoidance, persistence, and recurrence at stated quantitative
confidence --- unpacked by the logic's own proof theory into
Streett/tower verification conditions~\cite{abate24,abate26a} over
geolog. That proof theory
descends from Bernstein Logic, which survives in the finished design
as a named fragment. The discipline throughout is certificate-first:
untrusted, fast, freely sloppy searches \emph{propose}; a small exact
kernel \emph{disposes}. The certificates that discipline produces are
signed, content-addressed geolog commits, and three checking paths
consume those commits directly --- one set of bytes, no per-path
export format, three consumers built to fail differently:
\begin{itemize}
\item \texttt{geolog-alpha} is the production path, live today: the
  chase engine checks certificate rows incrementally as commits
  arrive, running QTSL's rules as it runs any theory. Its comparative
  advantage is throughput and integration; its acceleration lanes
  (GPU where the hardware offers them, SIMD elsewhere) are
  conformance-gated, so a lane may change cost, never a verdict. Its
  assurance is engineering-grade by deliberate policy: reference
  implementations, differential batteries, rebuild-invariance checks,
  and mutant corpora built to fool it, rather than mechanized
  theorems. That policy sets an assurance level; it does not offload
  soundness. For every certificate species beyond the Bernstein Logic
  fragment, this engine is today the only checker there is, and a
  certificate that fooled it would be a soundness failure of the
  toolkit. What separates it from \texttt{qtslab} is the size of the
  trusted base and the grade of the evidence --- mechanized down to
  the instruction set on one side, hardened against every attack we
  could construct on the other. On the Bernstein Logic fragment,
  whatever this engine derives can be re-checked, from the same
  bytes, by the two paths below.
\item \texttt{qtslab} is the certified path and the smallest trusted
  base of the three: a standalone binary whose soundness theorem ---
  machine-checked in HOL4 and grounded, through the CakeML verified
  compiler~\cite{kumar14cakeml}, in the semantics of the Arm
  instruction set itself ---
  says that if the binary prints ACCEPT on a claim followed by an
  arbitrary stream of bytes, the claim is semantically true; decode
  totality over arbitrary input is part of the proof obligation, so
  malformed bytes can end only in a refusal with a witness. It is
  assembling now, along a demonstrated route: decode totality is
  proven, and the same compiler-and-theorem recipe has already
  carried a complete predecessor checker to a machine-code soundness
  theorem on real hardware. \texttt{qtslab} only checks the Bernstein
  Logic fragment of QTSL, so it never sees a probability. The
  fragment --- sequents over the certificate's own declared hypernet
  (chapter~\ref{ch:qtsl-hypernets}), certified valid under every
  admissible valuation --- is all its input format carries today, and
  the boundary fails closed: a certificate reaching beyond it is
  refused, never silently accepted.
\item Lean is the independence path: the same bytes, accepted by the
  Lean kernel as a theorem via a reflection principle. The
  rule-by-rule soundness lemmas the calculus rests on are developed
  in Lean and load-bearing today; the checking route itself was
  expected to be too slow for machine-scale certificates, and was kept
  in the roster to measure that rather than assert it. It has been
  measured --- twice, across two certificate-format generations ---
  and the Experiments chapter carries the result: a scaling law and an
  admission gate for the compiled route, a kernel-mode envelope
  already exhausted below a kilobyte of certificate, and the
  expectation confirmed with its boundary stated. In exchange, the
  route ties the system to a proof kernel scrutinized by a large
  community that owes this toolkit nothing.
\end{itemize}
That boundary is \texttt{qtslab}'s, not the vocabulary's: the
certificate theory itself now encodes the full QTSL vocabulary ---
Systems, Specs, and the judgments connecting them --- as a
machine-checked schema, with strict well-formedness checks, a
termination certificate for its rule set, and fixture and mutant
batteries, all run under the \texttt{geolog-alpha} binary. That
theory set ships in this repository as data, battery beside it:
\path{theories/Qtsl.geolog}, extending the Bernstein Logic proof
tower of \path{theories/BernsteinLogic.geolog}. Producer-side
emission is live as well: the MTK producer seals certificate rows in
this vocabulary as single signed geolog commits, and the reproduction
spine already carries benchmark families from several of its
campaigns in that sealed form, mutant batteries and re-checking make
targets beside them. What remains open is
the growth of \texttt{qtslab}'s machine-checked coverage toward the
full logic. A certificate that
passes all three paths has been checked by artefacts that share no
implementation and no compiler and were built to fail differently; any one of them
alone is already a complete answer to ``why should I believe this?''
The Rule Catalogue chapter carries the calculus rule by rule in
literate form, each covered entry beside its mechanized statements
and signed check receipts, under a coverage dashboard recording how far
mechanized verification of soundness actually extends.

\textbf{MTK} is the producer stack: the machinery that makes
certificate production cheap enough to do at scale. It contributes the
numerically heavy proposers as GPU kernel rows --- Bernstein and
Chebyshev table construction, cell and simplex indexing, batched linear
programming, and eigensolves --- each returning witness-typed output that
the exact tier can dispose of. Its agent surface is specified and
arriving: a session protocol under which agents will drive tactic
selection over a stateful obligation workspace, with the strategy
intelligence supplied by the reader as the language model of their
choice. None of it is in the trusted base: not the proposers, not the
GPU code, not any language model driving the strategy --- all of those
may be arbitrarily wrong without compromising a certificate that
checks. Its cookbook chapter is agent-facing by design:
proof-production moves in the Polya shape, written in the reader's
native language.

The stance generalizes to this book's readers: anything impressive
ships with an assurance edifice small enough to audit end-to-end ---
here, the exact kernel plus whichever checking path the reader elects
--- and every load-bearing claim carries a check the reader can run.
Where the build deliberately departs from the programme's founding
documents, the departures are stated as narrowings, with their
reasons.

\section*{What is certified today}

The Experiments chapter (chapter~\ref{ch:experiments}) is this
volume's evidence: campaigns against benchmarks the toolkit does not
own, each result a concrete policy or ranking function together with a
machine-checkable certificate, every number re-derived from receipt
files at each preview. The headlines below are compressions of that
chapter, not claims made here: each resolves, by scrolling, to the
section carrying its numbers, its tables, and its stated limits
(what each claim species means, and how to read the tables, are
themselves sections there), with the clone-and-repro path beneath
that as the deeper tier. The chapter dates its own state, names what
is in flight, and grows at each preview.

\paragraph{A published verification table, re-done exactly.}
\looseness=-1 All
seventeen rows of a peer-reviewed neural-certificate results table
(LMCS 2026)~\cite{abate23} now carry strictly better bounds from exact-rational
certificates --- margins from $1.01\times$ on the near-tight rows to
$58$ orders of magnitude on the conservative ones --- and the exact
computation refuted one published ``sound'' bound. Each row's proof
decomposes into per-obligation certificate files --- $28{,}569$
across the campaign, most of them the $24{,}000$ lattice-piece
certificates behind the one continuous-state row --- and the full
corpus has replayed through three successive checker builds with zero
refusals.
(section~\ref{sec:exp-sweep}.)

\paragraph{The first certified rows for a standard RL suite.} On
Gymnasium's tabular tier: shipped lookup-table policies with
certificates of exact optimality --- among them the first published
exact value of Sutton--Barto blackjack under optimal play, a proof
that the suite's own Taxi ``solved'' threshold is unsatisfiable, and
belief-chain certificates for hidden-state variants (certifying the
implementable policy, never a clairvoyant one) --- each row checking
in microseconds to milliseconds. On the physics-engine tier, where
analytic claims would be a category error, anytime-valid statistical
certificates: LunarLander expected return $\ge 236.2$ (discrete) and
$\ge 281.1$ (continuous) at confidence $99\%$ over $10^6$ episodes
each --- to our knowledge the first certified statements of any kind
at this claim class for Box2D environments. And the continuous-state
analytic tier is now open: Pendulum carries a certified
\emph{worst-case} return bound, holding for every initial state with
no probability anywhere in the certificate.
(section~\ref{sec:exp-gym}.)

\paragraph{A hardware model-checking suite, completed with
certificates.} On the neural model-checking line's benchmark
(NeurIPS 2024/2025, eleven Verilog
designs)~\cite{giacobbe24,neuralmc-2025-invariant}: all $634$ tasks across
its three property classes certified, against $554$ for the
strongest published tool and $617$ for the per-task union of all
five published baselines --- every task shipping a small exact
certificate whose consuming lemma is machine-checked in Lean, where
every baseline ships a verdict. Replaying the whole suite --- every
certificate and every rewrite-chain witness --- takes under a second
on one core, against a published budget of five hours per task;
method invention to a full first sweep of the pure-liveness class
took about $75$ minutes, and each further class landed in a single
working session. (section~\ref{sec:exp-circuits}.)

\paragraph{An Atari console, compiled into a white-box environment.}
\looseness=-1
Pong and Freeway compile to synchronous circuits ($1{,}491$ flip-flops,
${\sim}21{,}000$ AND gates) with emulator-exact fidelity receipts ---
identical episode returns, $116$--$123$ of $128$ RAM bytes exact on
every frame, each divergence named --- so that certified claims about
the environment \emph{itself} become definable. The first certified
exact-replay rows arrived as this preview closed. Among them, on
Montezuma's Revenge (compiled the same way, its bank-switched
cartridge's mapper hardware included), are, to our knowledge, the
first certified returns for the suite's canonical hard-exploration
benchmark: the first key certified at every one of the $31$ no-op
init offsets, and a key-and-door trajectory certified at return
exactly $400$, bit-exact against the circuit's own transition
relation.
Certified policy rows are in flight. The scope condition is stated there
plainly: the circuit is examined, not learned --- system
identification is a separate problem the toolkit does not yet
address. (section~\ref{sec:exp-atari}.)

\paragraph{Semiconductor process models: certified stable, with
wall-clock dispatch optimal for up to seven lots on a simplified net.} On the
Intel Mini-Fab, the community's minimal re-entrant model, named
dispatch policies
carry exact drain certificates at up to $99.9\%$ of fluid capacity,
with the probe at capacity correctly failing; on a SEMATECH-scale
testbed, the certified load at the coarsest tractable aggregation is
$13.2\%$, and one walk of the refinement path built into the
certificate format lifts it to a certified $19.8\%$. Further
refinement is measured but not yet certified, and the numbers are
weak and shown anyway: the aggregation gap is a measured frontier,
not fine print. On the Mini-Fab's discrete net, dispatch itself is
now a certified quantity: joint policy+certificate synthesis returns
a supervisor with a certified $1.9\%$ throughput-floor gain over the
no-policy worst case; an explicit shipped policy attains the
wall-clock optimum over \emph{all} dispatch policies for up to seven
lots (two-sided, with LRAT proofs accepted by a CakeML-verified
checker, resting on two pen-and-paper lemmas), and attains the same
expression in replay at $N=84$ and at $N=8,\dots,12$ (with the
pending-lot count and unlock threshold substituted for the latter
replays). Optimality at those counts is open: the Lean
machine-contiguity floor is for an abstract schedule model. These
results are about a simplified net in which at most one lot per
process step may start, receive a unit of service, or finish per
tick. An uncertified black-box dispatcher runs inside a
certified switching envelope whose sealed certificate re-verifies
byte-identical after every run. One certificate covers lot counts
from $7$ to $10^{7}$ --- its container grows six bytes across the
sweep, and checking stays at ${\sim}2$\,ms.
(section~\ref{sec:exp-semicon}.)

\section*{Status of this volume}

As of this preview (the build stamped on the title page): the table
of contents is the book's
normative shape; this summary is a draft, to be regenerated from the
finished chapters; the QTSL rule
catalogue is generated in full from the mechanized rule sources; the
partial-completeness chapter is complete; the Experiments chapter is
a working draft; and the remaining chapters are full drafts (the
QTSL quick start and the proof-production tutorial carry listings
and transcripts from real, receipted runs). The material is unevenly
ripe by design, and each chapter states its own status, with draft
banners, absence boxes, and status notes marking what is settled,
what is staged, and what is still argued.
